% \section{结论}
\section{Conclusion}

% 本文研究了监督目标与推理协议如何共同影响科学写作评价模型。实验
% 结果表明，RS 虽然改善了教师 rationale 拟合和生成质量，但这一
% 优势并未转化为更好的评分性能。在序数任务中，RF 增加了严重错误，
% 因而在两种监督目标下均降低了 QWK；对错误 rationale 进行修正则
% 使大部分所选有害案例恢复正确。PaIR 和 RePaIR 进一步表明，配对
% 接口监督与 rationale 重生成能够缓解评分性能和接口适配之间的
% 冲突。这些结果说明，rationale 质量不能单独代表评分能力，监督
% 目标应与实际使用的推理协议联合设计和评估。

This paper examined how supervision objectives and inference protocols
jointly affect LLM-based scientific-writing evaluation. Although RS improved
teacher-rationale fitting and generated-rationale quality, these gains did not
translate into better scoring performance. On ordinal tasks, RF increased
severe errors and consequently reduced QWK under both supervision objectives,
whereas correcting erroneous rationales recovered most selected harmful cases.
PaIR and RePaIR further showed that paired-interface supervision and rationale
regeneration can mitigate the tension between scoring performance and interface
coverage. These findings indicate that rationale quality alone is not a
reliable proxy for scoring ability, and that supervision objectives should be
designed and evaluated together with the intended inference protocol.

% \section{局限性。}
\section{Limitations.}

% 本研究仅在 Qwen3-4B 和七个科学写作评价任务上采用固定 LoRA 配置
% 进行实验。此外，rationale 质量分析依赖外部 LLM 裁判，修正实验
% 也仅覆盖预先筛选的有害案例。因此，当前结果不能直接推广至其他
% 模型和任务，也不能证明生成的 rationale 忠实反映模型的内部推理。
% 未来工作应在更多模型、任务和完整样本上进行验证，并引入人工评价。

Our experiments are limited to Qwen3-4B, a fixed LoRA configuration, and seven
scientific-writing evaluation tasks. Moreover, rationale quality was assessed
by external LLM judges, and the correction experiment covered only selected
harmful cases. The findings may therefore not generalize to other models or
tasks, nor do they establish that generated rationales faithfully represent
internal reasoning. Future work should extend the evaluation to broader models,
tasks, and samples, with additional human assessment.